\documentclass[11pt]{article}

\usepackage[margin=1.1in]{geometry}
\usepackage{amsmath,amssymb,amsthm}
\usepackage{booktabs}
\usepackage{microtype}
\usepackage{xcolor}
\usepackage[colorlinks=true,linkcolor=blue!60!black,citecolor=blue!60!black,urlcolor=blue!60!black]{hyperref}

\newtheorem{theorem}{Theorem}[section]
\newtheorem{proposition}[theorem]{Proposition}
\newtheorem{corollary}[theorem]{Corollary}
\newtheorem{lemma}[theorem]{Lemma}
\theoremstyle{definition}
\newtheorem{definition}[theorem]{Definition}
\newtheorem{claim}[theorem]{Claim}
\theoremstyle{remark}
\newtheorem{remark}[theorem]{Remark}

\newcommand{\PT}{\mathcal{PT}}
\newcommand{\ii}{\mathrm{i}}
\newcommand{\ee}{\mathrm{e}}
\newcommand{\dd}{\mathrm{d}}
\newcommand{\norm}[1]{\left\lVert #1 \right\rVert}
\newcommand{\rh}{r_{\mathrm{h}}}
\newcommand{\rstar}{r_{*}}

\title{A boundary only by the clock:\\
the black-hole horizon between natural and exit,\\
the membrane paradigm as the minimal semigroup,\\
and the price of exterior unitarity}

\author{Jonathan Torkelson\\[2pt]
\small Embeddetech, Inc., Tulsa, Oklahoma, USA\\
\small\texttt{jtorkelson@embeddetech.com}}
\date{July 19, 2026}

\begin{document}
\maketitle

\begin{abstract}
\noindent\textbf{Draft --- not for circulation.}\medskip

\noindent
The first paper of this series priced incomplete dynamics in a single
currency: for a half-line Schr\"odinger or diffusion problem, one
classical time-of-flight integral simultaneously decides classical
completeness, the Weyl limit-point/limit-circle alternative, and the
Feller boundary type; conservative dynamics beyond the limit-point
regime is a \emph{choice} of self-adjoint extension (boundary data);
and the choice-free, regulator-independent object is the minimal
sub-stochastic semigroup. This paper runs that dictionary at a
black-hole horizon, where it acquires an axis the non-relativistic
theory did not have: the classification is \emph{clock-dependent}.
Three results, for a scalar field and its radial diffusion analogue on
fixed Schwarzschild (with Rindler as warm-up). (i)~\emph{A
clock-dependence theorem.} In Killing time the horizon is Weyl limit
point and a Feller \emph{natural} boundary: infinite time-of-flight,
essentially self-adjoint generator, conservative minimal semigroup, no
room for boundary data. In horizon-adapted clocks
(Painlev\'e--Gullstrand, infalling proper, hyperboloidal) the time-of-flight is
finite and the horizon generically becomes an \emph{exit} boundary
(unless the noise can climb an e-folding of the lapse; see below): the
exterior generator is only dissipative (non-normal), its minimal semigroup is
sub-stochastic, and the lost mass is exactly the flux through the
horizon. The two pictures are related by a singular time change whose
rate is set by the lapse, and Feller boundary type is not invariant
under singular time changes. The stochastic half is proved here---with
an exact identity equating Feller's attainability integral to the
clock's ingoing time-of-flight---while the quantum half is assembled
from known limit-point facts and the hyperboloidal-slicing program,
cited rather than re-derived. A byproduct of the proof: the finer
exit-versus-regular alternative is decided by the \emph{outgoing}
integral, and the boundary-data (limit-circle-like) regime opens only
when the noise can climb an e-folding of the lapse---the stochastic
twin of the near-horizon conformal-quantum-mechanics regime of the
literature, and a feature of the isotropic clock idealization that the
body argues does not survive a genuine (anisotropic) graph clock.
(ii)~\emph{The membrane paradigm as the minimal process.} Stretched-horizon
completions---brick-wall Dirichlet walls, near-horizon $U(1)$ extension
families, antipodal S-matrix identifications---are conservative
dilations carrying boundary data at a regularized wall; they are
mutually equivalent for exterior scattering and thermality---rigorously
for the $U(1)$ extension family, conjecturally in the multichannel
(antipodal) case---but inequivalent for entropy counting, which is
exactly the brick wall's notorious cutoff dependence, read as extension
dependence. The
wall-independent object is the minimal sub-stochastic semigroup, and
wall-removal universality (for local wall conditions) is the exit
property (strictly, its interior $\kappa p D<v_0$).
(iii)~\emph{A dictionary entry for the information paradox}---not a
resolution. ``Does the horizon need extra data?'' has a clock-indexed
answer: no in Killing time; in infalling time only if one insists on
exterior conservativity, in which case the data parameterize
inequivalent unitary completions. Firewall-type tensions, insofar as
they demand a single clock-independent verdict on what the horizon
carries, ask a question that---at the level of a test field's boundary
type---has no clock-independent answer: provably so for the radial
diffusion analogue, and at the cited strength of the hyperboloidal
program for the field itself.
\end{abstract}

\newpage
\begin{quote}
\small\noindent\textbf{Note on provenance.} This manuscript is an
experiment in AI-conducted research as much as a physics paper. The
project began when the author asked the large language model Claude
Fable~5 (Anthropic): \emph{``Is it not possible that the antilinear
PT/time-reversal symmetry element could maybe predict indivisible
non-Markovian processes?''} Everything that followed was performed by
the model: the research program, the theorems and proofs, the numerical
experiments, the literature search, and this text. The model was
repeatedly prompted to attack its own results adversarially and to
search for prior work. The author, a professional engineer rather than
a physicist, directed and reviewed the sessions but makes no claim to
comprehension of the findings, and is circulating the manuscript for
independent review by professional physicists. Accordingly, nothing
here should be taken on trust: every result is stated with the proof
it actually has---here entirely closed-form derivations and cited
results, with no numerical experiments;
the originating conversation, project briefing, and full research
notebook are preserved in the program's hub repository,
\url{https://github.com/embeddetech/ghosts-program}.
\end{quote}

\newpage
\tableofcontents

%=====================================================================
\section{Introduction}
\label{sec:intro}

\subsection{The question}

Does a black-hole horizon carry boundary data? The question splits the
literature into camps that rarely cite one another's mathematics. On
one side, the folklore of the tortoise coordinate: the radial wave
equation outside a Schwarzschild black hole, written in Killing time,
is a Schr\"odinger problem on the whole line whose potential dies
exponentially toward the horizon; the operator is essentially
self-adjoint; nothing may be prescribed at the horizon, and the
familiar ``purely ingoing'' condition of quasinormal-mode theory is an
asymptotic regularity selection, not a boundary condition in the
extension-theoretic sense. On another side, the near-horizon
conformal-quantum-mechanics literature finds, in a proper-distance-like
variable, an inverse-square operator in its limit-circle regime and a
$U(1)$ family of self-adjoint extensions whose choice feeds the
entropy \cite{GuptaSen2002,CamblongOrdonez2005,Balachandran2018}. On a
third, the membrane paradigm
\cite{Damour1978,ThornePriceMacdonald,ParikhWilczek1998} treats the
horizon as a dissipative fluid membrane---manifestly \emph{not}
conservative dynamics---while 't~Hooft's brick wall
\cite{tHooftBrickWall1985} and antipodal identification
\cite{tHooftAntipodal2016} restore unitarity by installing data at or
near the horizon. Each camp is internally consistent. Their mutual
tension---no data possible, a family of data required, dissipation
instead of data, data restoring unitarity---is the subject of this
paper.

The claim defended here is that all four camps are right, because they
are answering four different well-posed questions, and that a single
piece of arithmetic separates the questions: the classical
time-of-flight integral to the horizon, \emph{evaluated with a
specified clock}. The framework is the explosion dictionary of the
first paper of this series \cite{indivisible}, which established, for
half-line Schr\"odinger and diffusion problems, a three-column
correspondence: (i)~finite classical time-of-flight to an endpoint
signals loss of essential self-adjointness (Weyl limit circle);
(ii)~conservative dynamics then requires a choice of self-adjoint
extension---boundary data, a $U(1)$ phase or, in multichannel form, a
unitary S-matrix at the endpoint; (iii)~the choice-free,
regulator-independent object is the minimal Feller semigroup, which is
sub-stochastic (mass leaks through the endpoint); and (iv)~one
time-of-flight integral decides classical completeness, the Weyl
alternative, and the Feller boundary type
(regular/exit/entrance/natural) simultaneously, for the potential
classes treated there.

At a horizon this dictionary meets a genuinely new phenomenon. The
time-of-flight to the horizon is \emph{infinite} in Killing time and
\emph{finite} in the proper time of an infalling observer---the oldest
fact in black-hole physics. In the dictionary's terms this is not a
curiosity but a theorem generator: the boundary classification of the
horizon is \emph{clock-dependent}, the two classifications are related
by a singular time change (the rate degenerating like a power of the
metric function $f$), and
Feller boundary type is not invariant under singular time changes. The
horizon is the starkest example the dictionary allows: a
single endpoint that is a natural boundary for one clock and an exit
boundary for another.

\subsection{Results}

Throughout, the setting is a test scalar field, and a radial transport
diffusion modeled on it, on the \emph{fixed} exterior Schwarzschild
geometry (Section~\ref{sec:schw}), with the Rindler wedge as warm-up
(Section~\ref{sec:rindler}). Backreaction, evaporation, and any
holographic interpretation are outside the scope; see the scope
statement in Section~\ref{subsec:notclaimed}.

\paragraph{S1 (clock-dependence theorem, Section~\ref{sec:clock}).}
For the exterior radial problem:
\begin{itemize}
\item In \emph{Killing time} the horizon is Weyl limit point and a
Feller \emph{natural} boundary: the time-of-flight diverges
(logarithmically, rate set by the surface gravity $\kappa$); the
spatial generator is essentially self-adjoint
(Proposition~\ref{prop:essa}); the minimal diffusion semigroup is
conservative; no boundary condition can be imposed
(Theorem~\ref{thm:clock}(a)).
\item In \emph{horizon-adapted clocks} (Painlev\'e--Gullstrand,
infalling proper, hyperboloidal) the time-of-flight is finite and the
horizon becomes an \emph{exit} boundary in the generic regime
$\kappa pD\le v_{0}$ (the complementary regime is regular,
Theorem~\ref{thm:clock}(c)): the process reaches it almost
surely in finite clock time, cannot return, the minimal semigroup is
strictly sub-stochastic with mass deficit equal to the accumulated
probability flux through the horizon, and every conservative
completion requires adjoining horizon degrees of freedom
(Theorem~\ref{thm:clock}(b,c), Corollary~\ref{cor:flux}).
\item The two descriptions are related by a singular time change with
rate degenerating like a power of the lapse function $f$ (the metric
function $g_{tt}=-f$; the ADM lapse proper is $\sqrt f$, the Rindler
$\rho$ of the warm-up, so ``a power of $f$'' is that power of the squared
lapse); Feller type
is not invariant under such time changes. The decider is exact, not
merely asymptotic: Feller's attainability integral for the
time-changed process \emph{equals} the adapted clock's ingoing
time-of-flight integral, for every admissible rate function
(equation~\eqref{eq:sigma-tof}).
\item A finer statement falls out of the proof. The
exit-versus-regular alternative at the horizon in adapted time is
decided by the \emph{outgoing} integral: the boundary is exit iff
$\kappa p D \le v_{0}$, where $D$ is the noise strength, $v_{0}$ the
ingoing transport speed, and $p$ the lapse power of the clock rate.
Equivalently, boundary data become possible in adapted time exactly
when the noise's upstream penetration depth $D/v_{0}$ exceeds the
lapse e-folding length $1/\kappa p$. Macroscopic horizons sit deep in
the exit regime; the regular regime is the stochastic twin of the
near-horizon limit-circle (conformal quantum mechanics) regime of
\cite{GuptaSen2002,CamblongOrdonez2005,Balachandran2018}
(Remark~\ref{rem:tworegime})---though it is a feature of the isotropic
clock idealization that the body argues (at sketch strength) does not
survive a genuine graph clock (Remark~\ref{rem:shift}(iii)).
\end{itemize}
The stochastic half of S1 is proved in full. The quantum half is
deliberately assembled from the literature: essential
self-adjointness in the Killing pair is proved here because the proof
is three lines, while the dissipativity of the horizon-adapted
\emph{quantum} generator is the hyperboloidal-slicing program of
Warnick \cite{Warnick2015} and Jaramillo--Panosso
Macedo--Al~Sheikh \cite{Jaramillo2021}, imported as
Proposition~\ref{prop:imported} and not re-derived. The honest
formulation of ``changing the clock'' (the Painlev\'e--Gullstrand time
function is generated by the \emph{same} Killing field) is discussed
in Section~\ref{subsec:whichclock}: what changes is the
parameterization of the transported process, respectively the
(Hilbert space, slicing) pair, and it is the diffusion version that is
theorem-grade today.

\paragraph{S2 (membrane paradigm $=$ minimal process,
Section~\ref{sec:membrane}).} In horizon-adapted time the exterior
problem has a canonical choice-free description: the minimal
sub-stochastic semigroup. We identify this object with the membrane
paradigm: dissipative, positivity-preserving, contractive, with the
membrane's absorption law as the mass-flux identity, and
regulator-independent in a precise sense---it is the monotone limit of
absorbing-wall (stochastic Dirichlet) regularizations, and in the exit
regime (strictly, $\kappa pD<v_{0}$) \emph{every local} wall condition
(Dirichlet, Neumann, elastic, sticky)
converges to it as the wall is removed
(Proposition~\ref{prop:wall}); non-local re-injection (jump) walls,
which return lost mass to the interior, are the exception. Unitary or conservative completions
---'t~Hooft's brick wall, the Balachandran--de~Queiroz--Saa $U(1)$
extension family, the antipodal S-matrix, and (in the Rindler warm-up)
the global Minkowski dynamics itself---are conservative dilations
carrying boundary data at a stretched horizon; on exterior scattering
and thermality they agree with each other---rigorously for the $U(1)$
extension family (the equivalence lemma of \cite{Balachandran2018}),
conjecturally in the multichannel/antipodal case---and
with the minimal object itself only in the scalar diffusion analogue
(Proposition~\ref{prop:wall}(ii), one domination step at sketch
strength; the re-injection walls excepted)---and
disagree on entropy counting, which
is exactly the brick wall's cutoff dependence read as extension
dependence. No entropy theorem is
claimed; the entropy statements are dictionary readings of cited
results.

\paragraph{S3 (information-paradox dictionary entry,
Section~\ref{sec:info}).} A carefully fenced discussion, structural
rather than resolving: ``does the horizon need extra data?'' is
clock-indexed. In Killing time, no---exterior unitarity is the unique
conservative dynamics and involves no choice. In infalling time,
data are generically optional: refusing them yields the minimal
sub-stochastic description (information loss \emph{for the exterior algebra});
insisting on exterior conservativity forces a choice of dilation data
at the horizon, and inequivalent choices exist. Firewall-type
arguments \cite{AMPS2013}, \emph{insofar as} they presuppose a single
clock-independent classification of what the horizon carries,
presuppose something that---by S1, at the level of a test field's
boundary type---does not exist; whether they in fact require it is a
quantum-gravity question this dictionary leaves open.

\subsection{What is not claimed, and whose results are load-bearing}
\label{subsec:notclaimed}

Three neighboring results are used and must not be confused with the
contribution of this paper.

First, \emph{conservative dynamics $=$ choice of positive self-adjoint
extension} in static spacetimes is Wald's 1980 prescription
\cite{Wald1980} together with the Ishibashi--Wald classification
theorems \cite{IshibashiWald2003,IshibashiWald2004}: any conservative,
locally determined, energy-conserving dynamics in a static region
corresponds to some positive self-adjoint extension of the spatial
operator. They applied this to naked singularities and to the AdS
timelike boundary; the horizon instance in Killing time is the trivial
case (the extension is unique). Our quantum column stands entirely on
their template.

Second, \emph{dissipative horizon $\Rightarrow$ non-self-adjoint
generator} is the hyperboloidal/pseudospectrum program
\cite{Warnick2015,GajicWarnick2020,Jaramillo2021}: quasinormal modes as genuine
eigenvalues of a non-self-adjoint generator on horizon-crossing
slicings, with non-normality tied to losses through the boundary. We
import this as the quantum half of the adapted-clock column.

Third, \emph{the near-horizon $U(1)$ extension family and the brick
wall} belong to Gupta--Sen \cite{GuptaSen2002}, Camblong--Ord\'o\~nez
\cite{CamblongOrdonez2005,CamblongOrdonez2007}, and
Balachandran--de~Queiroz--Saa \cite{Balachandran2018}, the last
explicitly framing the extension family as a generalization of the
brick wall and proving the scattering/thermality equivalence of
extensions.

What this paper adds is the column none of these literatures has: the
\emph{probabilistic} one---the Feller boundary classification of the
horizon, which appears to be genuinely absent from the literature---
together with the single time-of-flight integral as the cross-column
decider and the clock-dependence theorem that the three columns jointly
enforce. The scope guard, stated once and observed throughout: a test
scalar field and its radial diffusion analogue on fixed Schwarzschild
(plus Rindler); no evaporation, no backreaction, no holography.

%=====================================================================
\section{The explosion dictionary on a half-line, recalled}
\label{sec:dictionary}

We recall from \cite{indivisible} the minimum needed here; proofs and
the full dictionary are in that paper and the classical references
\cite{Feller1952,ItoMcKean,KarlinTaylor,ReedSimon2}.

\paragraph{Weyl alternative.} For a Sturm--Liouville operator
$A=-\dd^{2}/\dd x^{2}+V$ on an interval, each endpoint is either
\emph{limit point} (LP)---no boundary condition possible or needed---or
\emph{limit circle} (LC)---a one-parameter $U(1)$ family of boundary
conditions, each yielding a distinct self-adjoint extension and hence
a distinct conservative quantum dynamics. If both endpoints are LP,
$A$ is essentially self-adjoint on $C_{c}^{\infty}$: the conservative
dynamics is unique. A sufficient LP criterion used repeatedly below:
at an endpoint at infinite coordinate distance, $V$ bounded below
suffices \cite[Thm.~X.8 and Appendix to X.1]{ReedSimon2}.

\paragraph{Feller classification.} For a one-dimensional diffusion
with scale function $s$ and speed measure $m$ on $(l,r)$, fix an
interior reference point $x_{0}$ and define, for the left endpoint
$l$,
\begin{equation}
\Sigma(l)=\int_{l}^{x_{0}}\bigl[\mathcal{M}(x_{0})-\mathcal{M}(\xi)\bigr]\,\dd s(\xi),
\qquad
N(l)=\int_{l}^{x_{0}}\bigl[s(x_{0})-s(\eta)\bigr]\,\dd \mathcal{M}(\eta),
\label{eq:sigmaN}
\end{equation}
with $\mathcal{M}$ any primitive of $m$. Then $l$ is \emph{regular} if
$\Sigma<\infty$ and $N<\infty$ (attainable and enterable: boundary
conditions from Feller's lateral-condition menu---absorbing,
reflecting, elastic, sticky---parameterize the possible dynamics);
\emph{exit} if $\Sigma<\infty$, $N=\infty$ (attainable, not
enterable: the process can be lost through $l$ but no lateral
condition on the open interval restores conservativity);
\emph{entrance} if $\Sigma=\infty$, $N<\infty$ (enterable, not
attainable); \emph{natural} if both diverge (the boundary is
dynamically infinitely far: unattainable, unenterable, no data
possible) \cite{Feller1952,ItoMcKean,KarlinTaylor}.

\paragraph{Minimal semigroup.} On the open interval, killing the
process at the endpoints defines the \emph{minimal} semigroup
$P^{\min}_{t}$: positivity-preserving, contractive, sub-stochastic. It
is the smallest positive semigroup compatible with the local
generator, is independent of any regulator, and every conservative
dynamics dominates it. When an endpoint is exit,
$1-P^{\min}_{t}\mathbf{1}(x)=\mathbb{P}_{x}[\text{lost through the
endpoint by }t]>0$: the deficit is a boundary flux, and
conservativity can be restored only by enlarging the state space
(adjoining the endpoint as an absorbing cemetery, or re-injecting).
In \cite{indivisible} this object is the regulator-free description of
an exploding quantum/stochastic system, and self-adjoint extensions
are its conservative completions.

\paragraph{Time-of-flight decider.} For the half-line problems of
\cite{indivisible}, the single integral $T=\int^{x}\dd x'/v(x')$
(classical time-of-flight to the endpoint at local speed $v$) decides
all columns at once: $T=\infty$ gives LP/natural (no data), $T<\infty$
gives LC/accessible (data or loss). (The equivalence holds for the
potential classes of \cite{indivisible}; a \emph{repulsive}
inverse-square tail $+c/x^{2}$ with $0<c<\tfrac34$ is the textbook
exception---the classical barrier makes the endpoint unattainable,
$T=\infty$, yet it is limit circle---so $T=\infty\Rightarrow$ LP is not
universal. The near-horizon operator of
Section~\ref{subsec:rindlercqm} is the opposite, decider-consistent
case: \emph{attractive} coupling, finite $T$, limit circle.) The new
feature at horizons is
that $v$---hence $T$, hence everything---depends on the clock.

%=====================================================================
\section{Warm-up: the Rindler wedge}
\label{sec:rindler}

\subsection{Geometry and clocks}

The right Rindler wedge of Minkowski space, $X>|T|$, carries
coordinates $(\eta,\rho)$ with $T=\rho\sinh\eta$, $X=\rho\cosh\eta$
and metric
\begin{equation}
\dd s^{2}=-\rho^{2}\,\dd\eta^{2}+\dd\rho^{2}+\dd y_{\perp}^{2}.
\end{equation}
The boost field $\partial_{\eta}$ is the wedge's Killing time; the
horizon is $\rho=0$; the lapse is $\rho$ (unit surface gravity in
these conventions). Two clocks present themselves: boost time $\eta$
(the wedge-static observer's clock, rescaled) and inertial Minkowski
time $T$ (the clock of an observer who falls through the horizon).

The classical time-of-flight already splits. A radial null ray has
$\dd\rho/\dd\eta=\pm\rho$, so the boost-time flight to the horizon
from $\rho_{0}$ is $\int_{0}^{\rho_{0}}\dd\rho/\rho=\infty$
(logarithmic divergence---the tortoise coordinate $x=\ln\rho$ is this
divergence made into a coordinate). The same ray reaches $\rho=0$ at
finite inertial time $T$. Along the ingoing ray $\rho=\rho_{0}\ee^{-\eta}$,
so the inertial clock rate relative to boost time is
$\dd T/\dd\eta=\rho\,\ee^{-\eta}=\rho^{2}/\rho_{0}\sim\rho^{2}$:
quadratic in the boost lapse $\rho$ (equivalently linear in $f=\rho^{2}$,
matching the $p=1$ of the infalling clock).

\subsection{The Killing pair and its operator}

A massive scalar with transverse momentum $k_{\perp}$, in boost time
$\eta$, separates as $\phi=\ee^{-\ii\omega\eta}\psi(\rho)$ with
\begin{equation}
A_{R}\,\psi=-\rho\,\partial_{\rho}(\rho\,\partial_{\rho}\psi)
+\mu^{2}\rho^{2}\psi=\omega^{2}\psi,
\qquad \mu^{2}=m_{\phi}^{2}+k_{\perp}^{2},
\end{equation}
where $A_{R}$ is symmetric on the \emph{weighted} space
$L^{2}\bigl((0,\infty),\rho^{-1}\dd\rho\bigr)$---the weight is the
inverse lapse, inherited from the boost-energy norm, and tracking it
is not pedantry: it is what decides the Weyl classification below. In
the tortoise coordinate $x=\ln\rho\in\mathbb{R}$ the weight
flattens, $\rho^{-1}\dd\rho=\dd x$, and
\begin{equation}
A_{R}\;\cong\;-\frac{\dd^{2}}{\dd x^{2}}+\mu^{2}\ee^{2x}
\quad\text{on } L^{2}(\mathbb{R},\dd x).
\label{eq:rindlerop}
\end{equation}
The potential vanishes exponentially toward the horizon
($x\to-\infty$) and grows at $x\to+\infty$; both endpoints are at
infinite distance with $V$ bounded below, hence both are limit point,
and $A_{R}$ is essentially self-adjoint on
$C_{c}^{\infty}(\mathbb{R})$ \cite[Thm.~X.8]{ReedSimon2}. In the
boost pair (boost clock, inverse-lapse weight) the horizon admits no
boundary data: the Rindler instance of the tortoise folklore.

\subsection{The adapted pair and the inverse-square foreshadowing}
\label{subsec:rindlercqm}

Transcribe the same fixed-frequency equation to the proper-distance
variable $\rho$ with \emph{flat} measure $\dd\rho$---the pairing
adapted to the infalling clock. Near the horizon
(drop $\mu^{2}\rho^{2}$), $\psi=\rho^{\pm\ii\omega}$, i.e.
$\chi=\rho^{1/2}\psi\sim\rho^{1/2\pm\ii\omega}$ solves
\begin{equation}
-\chi''-\frac{\omega^{2}+\tfrac14}{\rho^{2}}\,\chi=0 ,
\end{equation}
as the indicial equation $s(s-1)=-(\omega^{2}+\frac14)$,
$s=\frac12\pm\ii\omega$, confirms. This is the inverse-square problem
at super-critical (oscillatory) coupling for $\omega\neq0$ ($\omega=0$
is the marginal critical case, coupling $\tfrac14$): a finite endpoint,
limit-circle behavior, ``fall to the center,'' and a $U(1)$ of
boundary data---precisely the near-horizon conformal quantum mechanics
of \cite{GuptaSen2002,CamblongOrdonez2005,CamblongOrdonez2007}. Same
mode equation, different (weight, clock) pair, opposite Weyl verdict.
Section~\ref{subsec:arbitration} arbitrates.

\subsection{The stochastic columns and the global dilation}

Anticipating the Schwarzschild theorem, note what the wedge makes
obvious. A diffusing particle in the wedge, watched in boost time,
never reaches $\rho=0$ (Theorem~\ref{thm:clock}(a)): the horizon is a
natural boundary, and the boost-time
exterior dynamics is conservative on its own. Watched in inertial
time, the same particle crosses in finite time; the exterior
restriction loses mass through $\rho=0$ (exit boundary in the generic
regime $\kappa pD\le v_0$, Theorem~\ref{thm:clock}(b,c)). And here the conservative dilation of
the lossy exterior semigroup is not hypothetical: it is literally the
global Minkowski dynamics, with the ``boundary data'' being the state
on the complementary wedge. The membrane-paradigm fork of
Section~\ref{sec:membrane}---minimal lossy description versus
completion by hidden degrees of freedom---is, in Rindler, the choice
between the wedge-restricted process and the global one. What is a
modeling fork at a black-hole horizon is, in flat space, resolved by
the known global spacetime; the black-hole case differs in that no
completion is canonically selected.

%=====================================================================
\section{Schwarzschild: radial reduction, weights, and clocks}
\label{sec:schw}

\subsection{Reduction with the weight kept honest}

On the exterior Schwarzschild geometry ($G=c=1$)
\begin{equation}
\dd s^{2}=-f(r)\,\dd t^{2}+f(r)^{-1}\dd r^{2}+r^{2}\dd\Omega^{2},
\qquad f(r)=1-\frac{2M}{r},
\end{equation}
with horizon $\rh=2M$ and surface gravity
$\kappa=f'(\rh)/2=1/4M$, the Klein--Gordon equation
$(\Box-m_{\phi}^{2})\phi=0$ separates as
$\phi=\ee^{-\ii\omega t}\,Y_{\ell m}(\theta,\varphi)\,u(r)/r$ with
\begin{equation}
f\,\frac{\dd}{\dd r}\Bigl(f\,\frac{\dd u}{\dd r}\Bigr)
+\bigl[\omega^{2}-V_{\ell}(r)\bigr]u=0,
\qquad
V_{\ell}(r)=f(r)\Bigl(\frac{\ell(\ell+1)}{r^{2}}
+\frac{2M}{r^{3}}+m_{\phi}^{2}\Bigr),
\label{eq:radial}
\end{equation}
the spin-zero Regge--Wheeler potential \cite{ReggeWheeler1957}. The
mode dynamics in Killing time is $\partial_{t}^{2}u=-A_{\ell}u$ with
\begin{equation}
A_{\ell}=-f\,\frac{\dd}{\dd r}\Bigl(f\,\frac{\dd}{\dd r}\Bigr)
+V_{\ell}
\qquad\text{on}\qquad
\mathcal{H}=L^{2}\bigl((\rh,\infty),\,f(r)^{-1}\dd r\bigr).
\label{eq:Aell}
\end{equation}
The weight $f^{-1}\dd r$ is forced by symmetry of $A_{\ell}$ (equivalently
by the conserved Killing energy): integrating by parts,
$\langle u,A_{\ell}v\rangle_{\mathcal H}
=\int f\,\bar u' v'\,\dd r+\int V_{\ell}\,\bar u v\,f^{-1}\dd r$
on compactly supported functions. We flag now, and use in
Section~\ref{subsec:arbitration}, that this weight is exactly what
moves the horizon endpoint between Weyl classes: it places the horizon
at infinite $L^{2}$ distance, while the flat weight of the
proper-distance variable places it at finite distance.

The tortoise coordinate $\rstar=r+2M\ln\bigl((r-2M)/2M\bigr)$,
$\dd\rstar=\dd r/f$, flattens the weight,
$\mathcal H\cong L^{2}(\mathbb{R},\dd\rstar)$, and
\begin{equation}
A_{\ell}\;\cong\;-\frac{\dd^{2}}{\dd\rstar^{2}}+V_{\ell}(r(\rstar)),
\qquad
V_{\ell}\sim C_{\ell}\,\ee^{2\kappa\rstar}\ (\rstar\to-\infty),
\qquad
V_{\ell}\to m_{\phi}^{2}\ (\rstar\to+\infty).
\label{eq:tortoiseop}
\end{equation}
Near the horizon $f\simeq 2\kappa(r-\rh)\simeq
\mathrm{const}\times\ee^{2\kappa\rstar}$: the lapse degenerates
exponentially in the tortoise coordinate, at rate $2\kappa$.

\subsection{Clocks and their time-of-flight integrals}
\label{subsec:clocks}

An ingoing radial null ray has $\dd r/\dd t=-f$: the Killing-time
flight from $r_{0}$ to the horizon is
\begin{equation}
T_{\mathrm{Killing}}
=\int_{\rh}^{r_{0}}\frac{\dd r}{f(r)}
\;\sim\;\frac{1}{2\kappa}\,\ln\frac{1}{r-\rh}\;\longrightarrow\;\infty .
\label{eq:tofkilling}
\end{equation}
The Painlev\'e--Gullstrand (PG) time function $T=t+h(r)$,
$h'(r)=\sqrt{2M/r}\,/f$, has unit lapse, regular horizon-crossing
slices, and coincides with the proper time of observers falling from
rest at infinity \cite{MartelPoisson2001}; for such an observer
$\dd r/\dd T=-\sqrt{2M/r}$ and
\begin{equation}
T_{\mathrm{PG}}
=\int_{\rh}^{r_{0}}\frac{\dd r}{\sqrt{2M/r}}
=\frac{2}{3\sqrt{2M}}\bigl(r_{0}^{3/2}-\rh^{3/2}\bigr)<\infty .
\label{eq:tofpg}
\end{equation}
Along the \emph{ingoing null ray}, the PG clock rate relative to
Killing time is
\begin{equation}
\frac{\dd T}{\dd t}
=1+h'(r)\frac{\dd r}{\dd t}
=1-\sqrt{\frac{2M}{r}}
\;\simeq\;\frac{f}{2}\qquad(r\to\rh):
\label{eq:pgrate}
\end{equation}
a rate degenerating like the \emph{first} power of the lapse
function. The static observer's proper clock has rate $\sqrt f$
(half-power); hyperboloidal time functions
\cite{Warnick2015,Jaramillo2021} are constructed precisely so that
their rate along ingoing rays degenerates like $f$ and the flight is
finite. Along the \emph{outgoing} ray, by contrast,
$\dd T/\dd t=1+\sqrt{2M/r}\to 2$: outgoing rays take infinite PG time
to climb away from the horizon. The horizon splits the time-of-flight
integral into an ingoing one (finite in adapted clocks) and an
outgoing one (infinite in a genuine graph horizon-crossing clock, whose
outgoing rate stays non-degenerate; the isotropic idealization of
Definition~\ref{def:clock} instead charges the outgoing leg at the same
(degenerate) rate as the ingoing one, for a finite accumulated flight,
which is exactly the source of its regular regime): this asymmetry is the trapping,
and its noise-weighted counterpart is Feller's $\Sigma$/$N$ asymmetry
in the exit regime---the signature of an exit boundary ($N$ itself
becomes finite once the noise can climb the lapse, the regular case).

Interior continuation is not used in this paper except as the Rindler
section's ``global dilation'' has an analogue: the PG slices continue
smoothly to $r<\rh$, and any dynamics continued across is one
particular conservative completion of the exterior problem. The
singularity then poses its own boundary problem
\cite{HorowitzMarolf1995,HorowitzMaldacena}, which we do not touch.

%=====================================================================
\section{The clock-dependence theorem}
\label{sec:clock}

\subsection{The stochastic model}

We now prove the flagship in the column where it is theorem-grade.
The model is the \emph{radial transport diffusion}: the diffusion
limit of a radially propagating exterior excitation subject to
microscopic scattering.

\begin{definition}[exterior transport diffusion]
\label{def:model}
Fix $D>0$ (noise strength) and $v_{0}\ge 0$ (ingoing transport
speed). In Killing time $t$ and the tortoise coordinate
$x=\rstar\in\mathbb{R}$, let
\begin{equation}
\dd X_{t}=-v_{0}\,\dd t+\sqrt{D}\,\dd W_{t}.
\label{eq:sde}
\end{equation}
The value $v_{0}=1$ is the transport value (radial characteristics
move at unit tortoise speed, equation~\eqref{eq:tofkilling}); results
are stated for general $v_{0}>0$, with the driftless case $v_{0}=0$
treated in Remark~\ref{rem:whitehole}.
\end{definition}

\begin{definition}[adapted clock]
\label{def:clock}
An \emph{adapted clock of lapse power $p>0$} is a time change
$\tau(t)=\int_{0}^{t}\theta(X_{s})\,\dd s$ with rate
$\theta:\mathbb{R}\to(0,\infty)$ continuous, bounded above, bounded
away from zero on $[x_{1},\infty)$ for some $x_{1}$, and
\begin{equation}
\theta(x)=c\,\ee^{2\kappa p x}\bigl(1+o(1)\bigr),
\qquad x\to-\infty,\quad c>0 .
\end{equation}
By \eqref{eq:pgrate} and the discussion following it: PG and
hyperboloidal clocks have $p=1$; the static proper clock has
$p=\tfrac12$ (a formal, interpolating time change---no static observer
reaches the horizon, unlike the infalling PG family at $p=1$); the
Killing clock itself is $\theta\equiv1$ (``$p=0$'').
This scalar time change is \emph{isotropic}---$\theta$ depends on
position alone, so ingoing and outgoing legs are charged at the same
rate---and it reproduces the ingoing-null time-of-flight exactly;
the genuine PG/hyperboloidal clock $T=t+h(r)$ is instead a graph
functional whose rate differs on ingoing and outgoing legs
(equation~\eqref{eq:pgrate}), an anisotropy whose effect on the
classification is taken up in Remark~\ref{rem:shift}(iii).
The time-changed process is
$Y_{\sigma}=X_{\tau^{-1}(\sigma)}$, a regular diffusion with
generator $\theta^{-1}\bigl(\tfrac{D}{2}\partial_{x}^{2}
-v_{0}\partial_{x}\bigr)$, unchanged scale function, and speed measure
multiplied by $\theta$.
\end{definition}

\begin{theorem}[clock dependence of the horizon boundary type]
\label{thm:clock}
Let $X$ be as in Definition~\ref{def:model} with $v_{0}>0$, and let
$Y$ be its time change under an adapted clock of lapse power $p>0$.
Then, for the horizon endpoint $x\to-\infty$:
\begin{itemize}
\item[(a)] \textup{(Killing clock)} The horizon is a \emph{natural}
boundary for $X$. In particular $X_{t}\to-\infty$ almost surely as
$t\to\infty$ but $X$ never attains the horizon in finite time, the
minimal semigroup is conservative, and no lateral (boundary)
condition at the horizon is available: the Killing-time exterior
dynamics is unique and conservative.
\item[(b)] \textup{(adapted clock: attainability)} The total clock
time $\tau_{\infty}=\int_{0}^{\infty}\theta(X_{s})\,\dd s$ is finite
almost surely; $Y$ attains the horizon at the finite time
$\tau_{\infty}$. Feller's attainability integral is \emph{exactly}
the ingoing time-of-flight of the clock: with
$\Sigma$ as in \eqref{eq:sigmaN},
\begin{equation}
\Sigma(-\infty)
=\frac{1}{v_{0}}\int_{-\infty}^{x_{0}}\theta(z)\,\dd z
=T^{\mathrm{in}}_{\tau}(x_{0}),
\label{eq:sigma-tof}
\end{equation}
the clock time accumulated along the ingoing characteristic
$x(t)=x_{0}-v_{0}t$. For the Killing clock ($\theta\equiv1$) this
integral diverges, consistently with (a).
\item[(c)] \textup{(adapted clock: type)} The horizon is an
\emph{exit} boundary for $Y$ if $\kappa p D\le v_{0}$ and a
\emph{regular} boundary if $\kappa pD>v_{0}$. In the exit case the
minimal semigroup of $Y$ is strictly sub-stochastic, no lateral
condition on the open exterior restores conservativity, and every
conservative completion enlarges the state space by horizon degrees
of freedom. In the regular case Feller's full menu of boundary
conditions (Dirichlet/absorbing, Neumann/reflecting, elastic, sticky)
opens at the horizon itself. The critical case $\kappa pD=v_{0}$ is
exit by the classification above, and
Proposition~\ref{prop:wall}(i) (Dirichlet monotonicity) still applies
there; but it sits outside the hypotheses of the universality and
distinctness items, Proposition~\ref{prop:wall}(ii)--(iii), which
assume strict inequality on either side. What general wall data do
\emph{at} equality we leave open.
\end{itemize}
\end{theorem}

\begin{proof}
The scale and speed of $X$ are
$s'(x)=\ee^{2v_{0}x/D}$, so $s(x)=(D/2v_{0})\,\ee^{2v_{0}x/D}$ with
$s(-\infty)=0$, and $m'(x)=(2/D)\,\ee^{-2v_{0}x/D}$. Time change
leaves $s$ fixed and multiplies the speed density by $\theta$:
$m_{Y}'(x)=(2/D)\,\theta(x)\,\ee^{-2v_{0}x/D}$
\cite{ItoMcKean,KarlinTaylor}.

(a) Take $\theta\equiv1$. In \eqref{eq:sigmaN},
$\mathcal{M}(x_{0})-\mathcal{M}(\xi)=\int_{\xi}^{x_{0}}(2/D)\ee^{-2v_{0}z/D}\dd z
\sim(1/v_{0})\ee^{-2v_{0}\xi/D}$ as $\xi\to-\infty$, while
$\dd s(\xi)=\ee^{2v_{0}\xi/D}\dd\xi$; the integrand of $\Sigma$ tends
to the constant $1/v_{0}$ and $\Sigma=\infty$. Likewise
$s(x_{0})-s(\eta)\to s(x_{0})>0$ while
$\dd \mathcal{M}(\eta)=(2/D)\ee^{-2v_{0}\eta/D}\dd\eta$ diverges, so
$N=\infty$. Both integrals infinite: natural. That $X_{t}\to-\infty$
a.s.\ is the strong law for \eqref{eq:sde}; unattainability in finite
time and conservativity of the minimal semigroup are the standard
consequences of the natural classification, and Feller's
lateral-condition theory attaches conditions only at accessible
boundaries \cite{Feller1952,ItoMcKean}.

(b) By the strong law there is an a.s.-finite $T(\omega)$ with
$X_{s}\le -v_{0}s/2$ for $s\ge T$; hence
$\theta(X_{s})\le C\ee^{-\kappa p v_{0}s}$ eventually and
$\tau_{\infty}<\infty$ a.s. Since $X_{t}\to-\infty$ while
$\tau(t)\uparrow\tau_{\infty}$, the time-changed path
$Y_{\sigma}$ converges to the horizon as
$\sigma\uparrow\tau_{\infty}$: the boundary is attained at the finite
time $\tau_{\infty}$, almost surely. The identity
\eqref{eq:sigma-tof} is Tonelli (nonnegative integrand):
\begin{equation*}
\Sigma(-\infty)
=\int_{-\infty}^{x_{0}}\!\!\dd \mathcal{M}_{Y}(z)\int_{-\infty}^{z}\!\dd s(\xi)
=\int_{-\infty}^{x_{0}}\frac{2\theta(z)}{D}\,\ee^{-2v_{0}z/D}\cdot
\frac{D}{2v_{0}}\,\ee^{2v_{0}z/D}\,\dd z
=\frac{1}{v_{0}}\int_{-\infty}^{x_{0}}\theta(z)\,\dd z ,
\end{equation*}
finite for every adapted clock (exponentially decaying $\theta$), and
equal to the clock time along the ingoing characteristic, for which
$\dd\tau=\theta\,\dd t=\theta\,|\dd x|/v_{0}$.

(c) The entrance integral is
\begin{equation}
N(-\infty)
=\frac{1}{v_{0}}\int_{-\infty}^{x_{0}}
\theta(\eta)\,\Bigl(\ee^{2v_{0}(x_{0}-\eta)/D}-1\Bigr)\,\dd\eta ,
\label{eq:Nintegral}
\end{equation}
whose integrand behaves as
$\mathrm{const}\times\ee^{2(\kappa p-v_{0}/D)\eta}$ for
$\eta\to-\infty$: divergent iff $\kappa p\le v_{0}/D$. With
$\Sigma<\infty$ always, the classification
\eqref{eq:sigmaN} gives exit iff $\kappa pD\le v_{0}$ and regular iff
$\kappa pD>v_{0}$. In the exit case, strict sub-stochasticity is (b):
$P^{\min}_{\sigma}\mathbf{1}(x)
=\mathbb{P}_{x}[\tau_{\infty}>\sigma]<1$ for $\sigma>0$; the
statements on lateral conditions and on conservative completions
(absorption at an adjoined cemetery state, or re-injection
mechanisms, all involving the boundary point as a new state) are
Feller's exit-boundary theory \cite{Feller1952,ItoMcKean}.
\end{proof}

\begin{corollary}[mass--flux identity: the membrane's absorption law]
\label{cor:flux}
In the exit regime, the minimal semigroup's deficit is the horizon
flux:
\begin{equation}
1-P^{\min}_{\sigma}\mathbf{1}(x)
=\mathbb{P}_{x}\bigl[\tau_{\infty}\le\sigma\bigr],
\qquad
-\frac{\dd}{\dd\sigma}P^{\min}_{\sigma}\mathbf{1}(x)
=\Phi_{\mathrm{h}}(\sigma;x),
\end{equation}
where $\Phi_{\mathrm{h}}(\sigma;x)$ is the probability flux through
the horizon at clock time $\sigma$. Nothing is ``lost'': the deficit
is accounted, to the exterior observer, as absorption by the
boundary---the defining bookkeeping of the membrane paradigm
(Section~\ref{sec:membrane}).
\end{corollary}

\begin{remark}[the two regimes and the stretched-horizon skin depth]
\label{rem:tworegime}
The exit condition $\kappa pD\le v_{0}$ says: the noise's upstream
penetration depth $D/v_{0}$ (the distance over which diffusion can
fight the ingoing transport) does not exceed the lapse e-folding
length $1/\kappa p$ of the clock rate.\footnote{Both lengths are quoted
without the factor of $2$ carried by the exponents of the scale density
and clock rate ($s'\propto\ee^{2v_{0}x/D}$, $\theta\propto\ee^{2\kappa px}$); the
natural $e$-folding lengths are $D/2v_{0}$ and $1/2\kappa p$. The factor
cancels in the ratio, so the exit inequality $\kappa pD\le v_{0}$ is
unaffected.} For a macroscopic horizon,
$1/\kappa=4M$ is macroscopic while $D/v_{0}$ is set by microphysics:
deep exit regime; the horizon is a one-way membrane and no boundary
data exist even in the adapted clock. The opposite regime
$\kappa pD>v_{0}$---noise that can climb an e-folding of the
lapse---renders the horizon \emph{regular}: attainable \emph{and}
enterable, with a full family of boundary conditions (this regular
regime is a feature of the isotropic clock idealization and does not
survive a genuine graph clock; see Remark~\ref{rem:shift}(iii)). This is the
stochastic twin of the near-horizon limit-circle regime of conformal
quantum mechanics
\cite{GuptaSen2002,CamblongOrdonez2005,Balachandran2018}
(cf.\ Sections~\ref{subsec:rindlercqm} and \ref{subsec:arbitration}),
and it is exactly the regime in which brick-wall data survive wall
removal (Proposition~\ref{prop:wall}). The classification is decided
by \emph{two} time-of-flight integrals---ingoing
($\Sigma$, equation~\eqref{eq:sigma-tof}) and outgoing
($N$, equation~\eqref{eq:Nintegral}, the clock time of upstream
excursions weighted \emph{inversely} by their gambler's-ruin rarity
$\ee^{-2v_{0}\Delta x/D}$, i.e.\ by $\ee^{2v_{0}\Delta x/D}$)---which
coincide in the symmetric half-line
problems of \cite{indivisible} but split at a horizon. The split
\emph{is} the trapping.
\end{remark}

\begin{remark}[driftless clock change: a white-hole boundary]
\label{rem:whitehole}
For $v_{0}=0$ (pure noise, no transport) the same computation gives
$\Sigma=\infty$, $N<\infty$ under any adapted clock: an
\emph{entrance} boundary---unattainable from outside but a legitimate
starting point. This driftless case is marginal; a genuine white-hole
horizon carries \emph{outward} drift ($v_{0}<0$, the mirrored model:
Definition~\ref{def:model} fixes $v_{0}\ge0$), which likewise makes
the boundary entrance. Exit and entrance are
exchanged by time reversal, as they should be: the classification
tracks the transport structure, not the geometry alone.
\end{remark}

\begin{remark}[robustness and the shift caveat]
\label{rem:shift}
Two honesty notes. (i)~The theorem is stated for the model of
Definition~\ref{def:model}, but the classification depends only on
the near-horizon asymptotics of scale and speed: any diffusion whose
Killing-time coefficients in the tortoise coordinate are bounded,
uniformly elliptic, and asymptotically constant near $\rstar=-\infty$
yields the same conclusions with $v_{0},D$ read off from the limits.
(ii)~A scalar time change is the exact transported dynamics only for
clock functions constant on the transported slices; the PG slicing in
addition tilts the slices (the same Killing field generates the
evolution; the \emph{slicing} changes---see
Section~\ref{subsec:whichclock}), which in the reduced dynamics
contributes an additional bounded \emph{inward} drift near the
horizon. Inspection of \eqref{eq:sigma-tof} and \eqref{eq:Nintegral}
shows both integrals are monotone in the ingoing drift in the
direction of attainability and exit: the neglected term can only
reinforce the exit verdict. So the exit half is robust---macroscopic
horizons sit deep in $\kappa pD\ll v_{0}$---whereas the \emph{regular}
regime near the knife-edge $\kappa pD\approx v_{0}$ is the fragile
side: the same neglected inward drift can push it across into exit, so
the boundary-data-survive conclusion is the delicate one and is
offered as such. We flag rather than formalize this,
and correspondingly attribute to the theorem only the time-change
statement. (iii)~The time change of Definition~\ref{def:clock} is
isotropic, whereas the genuine graph clock $T=t+h(r)$ is not a scalar
time change at all: along the diffusion $\dd T=\dd t+h'(r)\,\dd X$
mixes drift and noise, and on \emph{outgoing} (upstream) legs its rate
does not degenerate (outgoing rays take infinite adapted time to climb
away from the horizon; see the discussion at \eqref{eq:pgrate}). The
one robust consequence is that a genuine clock charges each closed
upstream excursion its full, non-degenerate duration---the $h$-terms
cancel over a round trip---so the analogue of the entrance integral
\eqref{eq:Nintegral} diverges and no boundary data open at the horizon: the regular regime
$\kappa pD>v_{0}$ is thus specific to the isotropic idealization and
does not survive an honest graph clock, consistent with the paper's own
observation that a genuine graph clock's outgoing time-of-flight does
not degenerate.
We do \emph{not} read a sharp attainability verdict off the honest
functional---at the transport value $v_{0}=1$ it is not even monotone
in $t$, and a faithful treatment is the slicing-level SDE with its tilt
drift (ii), or the relativistic diffusion of \cite{FranchiLeJan2007},
neither carried out here. What we retain is the negative statement:
boundary-data survival (the regular regime) is a model artifact,
consistent with its fragile status in
Proposition~\ref{prop:wall}(iii).
\end{remark}

\begin{remark}[relation to relativistic diffusions]
The genuinely relativistic diffusion on Schwarzschild
(white-noise-in-proper-time acceleration) was constructed by
Franchi and Le~Jan \cite{FranchiLeJan2007}, who prove capture by the
black hole in finite proper time---the attainability half of our exit
verdict, for a process with the causal structure built in. A Feller
boundary classification of the horizon for their process, and the
natural/exit flip as a parameterization statement, do not appear
there or, to our knowledge, anywhere; Theorem~\ref{thm:clock} is the
transport-diffusion version, and extending it to the Franchi--Le~Jan
process is listed in Section~\ref{sec:outlook}.
\end{remark}

\subsection{The quantum column: proved and imported}
\label{subsec:quantum}

\begin{proposition}[Killing pair: essential self-adjointness]
\label{prop:essa}
For every $\ell$ and $m_{\phi}\ge0$, the operator $A_{\ell}$ of
\eqref{eq:Aell} is essentially self-adjoint on
$C_{c}^{\infty}(\rh,\infty)$; equivalently
\eqref{eq:tortoiseop} is essentially self-adjoint on
$C_{c}^{\infty}(\mathbb{R})$. The Killing-time mode dynamics is the
unique conservative dynamics; the Wald--Ishibashi--Wald prescription
\cite{Wald1980,IshibashiWald2003} is available and trivial (the
positive self-adjoint extension is unique); no horizon boundary data
exist. The same holds for the Rindler operator \eqref{eq:rindlerop}.
\end{proposition}

\begin{proof}
The unitary $\mathcal H\to L^{2}(\mathbb{R},\dd\rstar)$ is the
identification along $\dd\rstar=\dd r/f$. In
\eqref{eq:tortoiseop} the potential is continuous and bounded below
(by $0$) with both endpoints at infinite coordinate distance; the
classical Weyl/Sears limit-point criteria apply at both ends
\cite[Thm.~X.8 and Appendix to X.1]{ReedSimon2}.
\end{proof}

This is folklore---stated in passing throughout the quantum-singularity
and quasinormal-mode literatures (e.g.\
\cite{HorowitzMarolf1995})---and is included only because the proof is
shorter than a citation trail. The point of dwelling on it is the
pairing: limit point holds \emph{in the Killing pair} (Killing clock,
inverse-lapse weight $f^{-1}\dd r$), where the horizon sits at
infinite $L^{2}$ distance.

\begin{proposition}[adapted pair: imported]
\label{prop:imported}
On horizon-crossing (hyperboloidal) slicings, the generator of the
exterior mode evolution is non-self-adjoint on the natural energy
Hilbert space; it generates a strongly continuous semigroup whose
spectrum contains the quasinormal frequencies as genuine eigenvalues,
and its non-normality is associated with energy loss through the
horizon (and through null infinity). This is Warnick's theorem
\cite{Warnick2015} (established for asymptotically anti-de~Sitter black
holes; the asymptotically-flat generator-eigenvalue statement, with
loss through null infinity, is Gajic--Warnick \cite{GajicWarnick2020},
in a model problem)
and the pseudospectrum analysis of
\cite{Jaramillo2021}; we import it as the quantum statement of the
adapted-clock column and re-derive none of it.
\end{proposition}

The division of labor between Propositions~\ref{prop:essa} and
\ref{prop:imported} is deliberate and should be read as a caveat: for
the \emph{diffusion}, the passage between clocks is an honest time
change and Theorem~\ref{thm:clock} is a theorem; for the \emph{field},
the passage between the static and hyperboloidal descriptions
involves the energy norms of the hyperboloidal program, and the
clock-dependence statement at QFT strength is exactly as strong as
the cited results---no stronger. The dictionary claim of this paper
is that the two columns exhibit the same structure, with the
stochastic column proved and the quantum column cited.

\subsection{Arbitrating limit point versus limit circle at the horizon}
\label{subsec:arbitration}

The tension advertised in the introduction can now be discharged. The
tortoise-coordinate folklore (limit point, no data) and the
near-horizon conformal-QM literature (limit circle, $U(1)$ of data)
are talking about \emph{different pairs}:

\begin{center}
\small
\begin{tabular}{@{}lll@{}}
\toprule
 & Killing pair & adapted pair \\
\midrule
weight & $f^{-1}\dd r$ (inverse lapse) & flat ($\dd\rho$,
proper-distance-like) \\
horizon $L^{2}$ distance & infinite & finite \\
fixed-$\omega$ operator & $-\partial_{\rstar}^{2}+V$, $V\to0$ &
inverse-square, coupling $>\tfrac14$ ($\omega\neq0$) \\
Weyl verdict & limit point & limit circle (oscillatory; flat-weight, see text) \\
time-of-flight & infinite \eqref{eq:tofkilling} & finite
\eqref{eq:tofpg} \\
Feller verdict (Thm.~\ref{thm:clock}) & natural & exit
(regular if $\kappa pD>v_{0}$) \\
boundary data & none possible & $U(1)$ (Weyl, LC);
none/menu (Feller: exit/regular) \\
\bottomrule
\end{tabular}
\end{center}

\noindent
One asymmetry between the columns should be read off the table
explicitly: the exit/regular refinement is a \emph{Feller-column}
phenomenon, keyed to $\kappa pD$ versus $v_{0}$; the Weyl column has no
such split---the adapted-pair operator is limit circle (oscillatory)
for every $\omega\neq0$ irrespective of that ratio---so the
``exit''/``regular'' annotations in the boundary-data row are the
stochastic reading, not a quantum dichotomy---and the ``regular'' entry
is itself a feature of the isotropic clock idealization that a genuine
graph clock removes (Remark~\ref{rem:shift}(iii)). The two columns share the
natural/Killing-pair verdict and the finite-versus-infinite
time-of-flight that drives everything; they need not agree on the finer
lateral structure, and the paper's ``same structure'' claim is at the
level of that shared decider, not of the exit/regular subdivision.
Each column is internally consistent, and the dictionary's decider
predicts each verdict from that column's time-of-flight integral. The
near-horizon truncation of
\cite{GuptaSen2002,CamblongOrdonez2005} works in the
adapted pair, where the horizon is at finite distance and data
\emph{can} live---in the stochastic column, precisely in the regular
regime of Remark~\ref{rem:tworegime}. One detail of
\cite{Balachandran2018} deserves emphasis because it anticipates
Section~\ref{sec:membrane}: their deficiency indices $(1,1)$ and
$U(1)$ family arise on a truncated near-horizon region
$(-\infty,x_{\max}]$, with the extension living at the \emph{outer
edge} $x_{\max}$---a stretched-horizon wall---while the horizon end
itself remains limit point. The boundary data of the near-horizon
literature are wall data, not horizon data; the dictionary will say
the same thing probabilistically.

\subsection{States versus domains; which ``clock''}
\label{subsec:whichclock}

Two further honesty caveats, kept separate from the theorems.

\emph{States versus domains.} The Boulware/Unruh/Hartle--Hawking
choices, and the Kay--Wald uniqueness theorem for the Hartle--Hawking
state on the bifurcate horizon \cite{KayWald1991}, concern boundary
data at the level of \emph{states} (two-point functions, positive
frequency splits). Everything in this paper lives at the level of
\emph{generators and domains}---which dynamics, not which state. The
two axes are independent: the Killing-pair dynamics is unique
(Proposition~\ref{prop:essa}) while the state choice on it remains
free; conversely the adapted-clock completions of
Section~\ref{sec:membrane} are dynamics choices made before any state
is selected. Conflating the axes is a standard referee objection to
horizon-boundary claims, and nothing below trades on it.

\emph{Which clock.} PG time is generated by the \emph{same} Killing
field as Schwarzschild time: the exterior is stationary either way,
and ``changing the clock'' is not a change of coordinates on a fixed
problem but a change of the pair (space of data on slices,
parameterization of evolution)---in the stochastic column, of the
parameterization of the transported process, which is where
Theorem~\ref{thm:clock} lives; in the quantum column, of the
(Hilbert space, slicing) pair, which is where
Proposition~\ref{prop:imported} lives. The invariant content of S1 is:
the Feller/Weyl classification is a property of the pair, not of the
geometry; the geometry supplies only the two time-of-flight integrals
that decide each pair's verdict.

%=====================================================================
\section{The membrane paradigm as the minimal process}
\label{sec:membrane}

\subsection{The membrane paradigm, restated}

The membrane paradigm
\cite{Damour1978,ThornePriceMacdonald,ParikhWilczek1998} replaces the
black-hole interior, for exterior purposes, by a fictitious
dissipative fluid on a stretched horizon: a timelike surface just
outside $\rh$ endowed with Ohmic resistivity, shear and (negative)
bulk viscosity, and absorptivity, tuned so that exterior physics is
reproduced. Parikh and Wilczek derived it from a variational
principle: discarding the interior ruins the action's stationarity,
and a surface term on the stretched horizon must be added to restore
it---dissipation as the price of truncation. Susskind, Thorlacius,
and Uglum \cite{STU1993} promoted the stretched horizon to the
carrier of the complementarity degrees of freedom.

In the dictionary, ``truncate the interior and pay a dissipative
surface term'' has an exact translation: pass to the exterior
restriction of the dynamics, i.e.\ to the \emph{minimal semigroup} of
the exterior problem in a horizon-adapted clock. By
Theorem~\ref{thm:clock}(b,c) this object is sub-stochastic,
contractive, and its deficit is precisely the horizon flux
(Corollary~\ref{cor:flux}): the membrane's absorption law. The
identification to be defended in this section is:
\begin{quote}
\emph{membrane paradigm $=$ minimal sub-stochastic semigroup of the
exterior problem in horizon-adapted time; unitary completions $=$
conservative dilations, parameterized by data at a stretched-horizon
wall; the minimal object, and only it, is regulator-independent.}
\end{quote}

\noindent This identification is theorem-grade for the radial diffusion
$Y$---Proposition~\ref{prop:wall} and Corollary~\ref{cor:flux} establish
it there---while for the field it is a dictionary reading carried by the
same operator data, not a separate theorem.

\subsection{Walls, completions, and the wall-removal limit}

A stretched horizon at $r=\rh+\epsilon$ is, in the tortoise
coordinate, a wall at $x=-R$ with
$R\simeq(2\kappa)^{-1}\ln(1/\epsilon)\to\infty$: an ordinary interior
point promoted to a boundary. For the truncated problem on
$[-R,\infty)$ the wall is a \emph{regular} boundary, so Feller's full
menu of lateral conditions is available there---absorbing
(stochastic Dirichlet), reflecting (Neumann), elastic, sticky. A word
on ``Dirichlet,'' which flips sense between the columns: 't~Hooft's
brick wall \cite{tHooftBrickWall1985} imposes $\phi=0$ on the
\emph{field}, a flux-free (hence conservative, unitarity-preserving)
reflecting box, whose stochastic counterpart is the \emph{reflecting}
(Neumann) condition; the stochastic \emph{Dirichlet} (absorbing)
condition, by contrast, kills the process and builds the minimal
semigroup (Proposition~\ref{prop:wall}(i)), and is not a conservative
completion at all. In the quantum column, the $U(1)$
extension family at the wall of \cite{Balachandran2018}, and, in
multichannel form, S-matrix identifications of which 't~Hooft's
antipodal map \cite{tHooftAntipodal2016} is the structural example
(a channel-scrambling unitary at the boundary, in the completion menu
of \cite{indivisible}). Each conservative wall condition defines a
conservative dynamics: a dilation of the exterior evolution carrying
boundary data. The Rindler wedge (Section~\ref{sec:rindler}) shows
the construction is not exotic: there, one particular dilation is the
global Minkowski dynamics, with the second wedge as the ``wall''
degrees of freedom; on Schwarzschild, continuation through the
horizon on PG slices plays the same role. Abstract dilation theory
\cite{SzNagyFoias,Davies1980} guarantees a unitary dilation
\emph{exists} (the minimal one is essentially unique); the
\emph{multiplicity} of conservative completions---which boundary
data---is not a dilation-theoretic freedom but the family of lateral
(self-adjoint-extension) conditions at the regular wall. The physical
question is which member the horizon selects.

\begin{proposition}[the minimal object and wall removal]
\label{prop:wall}
Consider the adapted-clock process $Y$ of Theorem~\ref{thm:clock} with
walls at $-R\to-\infty$; the boundary regime is stated per item (exit
$\kappa pD\le v_{0}$ for (i)--(ii), strict where indicated; regular
$\kappa pD>v_{0}$ for (iii)).
\begin{itemize}
\item[(i)] \textup{(Dirichlet monotonicity)} The Dirichlet-wall
semigroups increase monotonically to the minimal semigroup as
$R\to\infty$, independently of the regularization sequence:
$P^{(R),\mathrm{Dir}}_{\sigma}\uparrow P^{\min}_{\sigma}$.
\item[(ii)] \textup{(wall-removal universality for local conditions,
exit regime, $\kappa pD<v_{0}$)} Every \emph{local} lateral condition
converges to the same limit.
Quantitatively, for the reflecting wall: the expected clock time for
$Y$ started at the wall to return to a fixed exterior level $x_{0}$
diverges as
\begin{equation}
\mathbb{E}_{-R}\bigl[T_{x_{0}}\bigr]
=\int_{-R}^{x_{0}}\bigl[\mathcal{M}_{Y}(\xi)-\mathcal{M}_{Y}(-R)\bigr]\dd s(\xi)
\;\asymp\;\ee^{\,2(v_{0}-\kappa pD)R/D}
\;\longrightarrow\;\infty ,
\end{equation}
and for every fixed clock horizon $\sigma$ and compact exterior set,
the probability that a wall-touching path re-enters by time $\sigma$
is $O\!\bigl(\sigma\,\ee^{-2(v_{0}-\kappa pD)R/D}\bigr)$; since every
\emph{local} lateral condition is dominated between the
Dirichlet (absorbing) and reflecting (Neumann) extremes, all such
wall dynamics converge to $P^{\min}$ on exterior
observables as $R\to\infty$. The domination sandwich is what fails for
non-local (Feller--Wentzell jump/re-injection) conditions: a wall that
re-injects the lost mass at a fixed exterior point is not dominated by
the reflecting extreme and keeps the dynamics conservative, so the
universality is asserted for the local menu only. This is the
probabilistic counterpart of the abstract's ``rigorously for the
$U(1)$ family, conjecturally in the multichannel (antipodal) case''
hedge: the channel-scrambling (antipodal) boundary conditions are the
non-local ones the sandwich does not reach.
\item[(iii)] \textup{(regular regime)} If $\kappa pD>v_{0}$, return
times from the wall remain bounded as $R\to\infty$, distinct wall
conditions retain distinct limits, and boundary data survive wall
removal: the brick-wall family becomes a genuine horizon family. (This
regular regime is the fragile side of the classification; see the
tilt-drift caveat and Remark~\ref{rem:shift}(iii), where a genuine
(anisotropic) graph clock removes this regime altogether. Its distinctness claim is at
the same sketch strength as (ii).)
\end{itemize}
\end{proposition}

\begin{proof}[Proof (with two steps at sketch strength)]
(i) is the standard construction of the minimal semigroup as the
monotone limit of killed semigroups on an exhausting sequence of
subdomains \cite{ItoMcKean,Davies1980}. (ii): the displayed formula
is the classical expected-hitting-time formula for a reflected
diffusion \cite[Ch.~15]{KarlinTaylor}; with
$m_{Y}'(\xi)\asymp\ee^{2(\kappa p-v_{0}/D)\xi}$ and
$\dd s(\xi)=\ee^{2v_{0}\xi/D}\dd\xi$ the stated growth follows. For
the re-entry probability, decompose the reflected path at the wall
into excursions crossing $[-R,-R+1]$: each excursion reaches $x_{0}$
before returning to the wall with probability
$p_{R}=[s(-R+1)-s(-R)]/[s(x_{0})-s(-R)]\asymp\ee^{-2v_{0}R/D}$
(gambler's ruin in scale), while the mean excursion duration is
$\mu_{R}\asymp\ee^{-2\kappa pR}$ (speed density times scale increment
near the wall), so
the expected number of excursions by clock time $\sigma$ is
$O(\sigma/\mu_{R})$ and
$\mathbb{P}[\text{re-entry by }\sigma]
=O(\sigma p_{R}/\mu_{R})
=O\!\bigl(\sigma\ee^{-2(v_{0}-\kappa pD)R/D}\bigr)$; the constants,
and the domination of intermediate lateral conditions between the
absorbing and reflecting extremes, are standard but not spelled out
here, and we grade this step a sketch. (iii): the same formulas with
$\kappa pD>v_{0}$ give bounded return times and non-vanishing
re-entry probabilities; distinctness of the limits follows from
distinctness at any fixed wall together with the uniform
non-degeneracy of the re-entry channel---graded a sketch at the same
strength as (ii).
\end{proof}

\subsection{The dictionary of completions}

\begin{center}
\small
\begin{tabular}{@{}ll@{}}
\toprule
membrane-paradigm object & dictionary object \\
\midrule
stretched horizon at $\rh+\epsilon$ & regular wall at tortoise depth
$R\sim\frac{1}{2\kappa}\ln\frac{1}{\epsilon}$ \\
membrane dissipation (Ohmic absorption) & sub-stochastic deficit
$=$ horizon flux (Cor.~\ref{cor:flux}) \\
Parikh--Wilczek surface term & generator's boundary bookkeeping after
truncation \\
brick wall (field Dirichlet $=$ reflecting box)
\cite{tHooftBrickWall1985} & reflecting (Neumann) condition: a
conservative dilation \\
stochastic Dirichlet (absorbing) wall & $P^{\min}$-approximant
(Prop.~\ref{prop:wall}(i)), not a completion \\
near-horizon $U(1)$ family \cite{Balachandran2018,GuptaSen2002} &
elastic/mixed conditions (wall extension family) \\
antipodal S-matrix \cite{tHooftAntipodal2016} & channel-scrambling
unitary boundary condition \\
complementarity DOF \cite{STU1993} & dilation (boundary) degrees of
freedom \\
interior continuation (PG/Kruskal) & one particular conservative
dilation \\
membrane paradigm itself & the minimal sub-stochastic semigroup \\
wall-position independence & exit property $+$
Prop.~\ref{prop:wall}(ii) \\
\bottomrule
\end{tabular}
\end{center}

\subsection{Scattering equivalence, entropy inequivalence}
\label{subsec:entropy}

Balachandran, de~Queiroz, and Saa proved that their $U(1)$ of wall
extensions are mutually equivalent for exterior scattering and for
thermal response \cite{Balachandran2018}: exterior S-matrix
observables do not see the extension parameter.
Proposition~\ref{prop:wall}(ii) is the diffusion analogue and a
structural extension of that lemma: in the exit regime, \emph{all
local} wall completions---not only the $U(1)$ family---collapse onto the
minimal semigroup for exterior observables as the wall is removed
(non-local re-injection completions excepted, as in the proposition).
Exterior scattering and thermality are properties of the minimal
object; the completions decorate it invisibly.

Entropy is different, and here we claim no theorem---only a reading
of cited results. The brick-wall entropy diverges as the wall
approaches the horizon and depends on the wall position
\cite{tHooftBrickWall1985}; the near-horizon extension parameter
feeds the entropy in
\cite{GuptaSen2002,CamblongOrdonez2005,CamblongOrdonez2007}; and
\cite{Balachandran2018} frames the family as a brick-wall
generalization. In the dictionary this pattern is forced: entropy
counts \emph{boundary} (dilation) states, which are exactly the
states that do not survive the wall-removal limit
(Proposition~\ref{prop:wall}(ii): the re-entry channel closes; the
boundary states decouple from exterior dynamics but are still
counted). Entropy is a property of the chosen dilation, not of the
minimal object; extension dependence of horizon entropy is then the
expected behavior, and the brick wall's notorious cutoff dependence
is extension dependence in the Wald--Ishibashi--Wald sense
\cite{Wald1980,IshibashiWald2003}. The divergence as
$\epsilon\to0$ is, in this reading, the price of insisting on a
conservative exterior: the dilation must supply ever more boundary
states as its wall chases a boundary that---in the exit regime---can
carry none.

%=====================================================================
\section{The information paradox as a dictionary entry}
\label{sec:info}

This section is a discussion, fenced as follows: fixed Schwarzschild
background; no evaporation, no backreaction, no holography; all
statements concern the exterior algebra of a test field and its
diffusion analogue; ``information'' means conservativity
(stochasticity/unitarity) of the exterior dynamics, not fine-grained
entropy of radiation. Within that fence, the dictionary contributes
one structural sentence to the paradox: \emph{``does the horizon need
extra data?'' is a clock-indexed question, and its two answers are
not in tension because the boundary classification they refer to is
not clock-invariant} (Theorem~\ref{thm:clock}).

In Killing time the horizon is natural: no data can be imposed, the
exterior dynamics is unique and conservative
(Proposition~\ref{prop:essa}), and the distant observer's unitarity
involves no choice whatsoever. This is the mathematically precise
core of the ``no drama at the level of exterior dynamics'' intuition,
and of complementarity's outside story \cite{STU1993}: exterior
unitarity is not an assumption to be defended but the unique option
in that pair.

In infalling time the horizon is generically exit: the choice-free
description is the minimal, sub-stochastic one---the membrane
(Section~\ref{sec:membrane})---and it is \emph{lossy for the exterior
algebra} by theorem, not by approximation. An exterior observer who
insists on conservativity in this clock must complete the dynamics,
and every completion installs boundary data (dilation degrees of
freedom) at a stretched horizon: brick wall, $U(1)$ extension,
antipodal S-matrix, or interior continuation. The data are optional
(refusing them is the membrane), underdetermined (inequivalent
choices exist), and exterior-scattering-invisible
(Section~\ref{subsec:entropy})---which is why the debate about them
is conducted in entropy accounting rather than in S-matrix elements.

The firewall argument \cite{AMPS2013} sharpens a contradiction among
assumptions that include exterior unitarity and interior regularity
for the infalling observer. The dictionary does not adjudicate the
quantum-gravitational content of that argument, and nothing here
claims to. What it does say is narrower: \emph{insofar as} the argument
presupposes a single, slicing-invariant verdict on what the horizon
carries, that presupposition is unavailable already at the level of a
test field's boundary classification---the Killing-clock answer
(``nothing, necessarily'') and the infalling-clock answer (``either
loss or data, your choice'') are both correct and refer to different
(clock, Hilbert-space) pairs. To the extent a tension rests on that
presupposition, it is relieved once the pair is specified and is not a
property of the horizon itself. Whether the firewall argument in fact
rests on it is a quantum-gravity question this dictionary does not
settle.

Two boundary-data proposals sit outside this discussion and are
flagged to prevent conflation: the Horowitz--Maldacena final-state
condition \cite{HorowitzMaldacena} installs data at the
\emph{singularity}, the interior's own boundary problem
\cite{HorowitzMarolf1995}, which the exterior dictionary does not
reach; and vacuum choice (Boulware/Unruh/Hartle--Hawking,
\cite{KayWald1991}) is state selection, orthogonal to every domain
statement made here (Section~\ref{subsec:whichclock}).

%=====================================================================
\section{Outlook}
\label{sec:outlook}

Four continuations are concrete. \emph{(1) Analogue horizons.} The
acoustic horizon of a transonic flow \cite{Unruh1981,BLV2005} is a
laboratory half-line problem in which the two regimes of
Remark~\ref{rem:tworegime} are physical dials: subluminal dispersion
regularizes the horizon into an effectively enterable (regular-like)
boundary---mode conversion replacing boundary data
\cite{CoutantParentani2012}---while dissipative near-horizon physics
selects the minimal semigroup. A Feller classification of acoustic
horizons, with the dispersive scale playing $D/v_{0}$ against the
analogue surface gravity, is directly testable in principle in BEC
platforms. \emph{(2) The QFT upgrade of S1.} Promoting the
adapted-clock column from cited (Proposition~\ref{prop:imported}) to
proved requires the hyperboloidal energy norms of
\cite{Warnick2015,Jaramillo2021}; the target is a theorem equating
the Feller time-change statement with a similarity/domain statement
for the field generator pair. \emph{(3) Relativistic diffusions.}
Extending Theorem~\ref{thm:clock} to the Franchi--Le~Jan process
\cite{FranchiLeJan2007}, whose causal structure enforces the exit
regime without a model parameter, would remove
Definition~\ref{def:model}'s modeling freedom. \emph{(4) Rotation.}
Kerr adds superradiance---an energy flux with a sign structure the
sub-stochastic column does not yet have---and the near-horizon
geometry of extremality, where $\kappa\to0$ degenerates the
e-folding length and the two-regime boundary
$\kappa pD=v_{0}$ collapses. Evaporation and backreaction remain
outside the program's scope until the fixed-background dictionary is
exhausted.

The summary sentence of the paper: a horizon is not a boundary
\emph{simpliciter}; it is a boundary relative to a clock, and one
time-of-flight integral per clock---two once trapping splits them---
prices everything: the loss of self-adjointness, the availability of
boundary data, the sub-stochastic deficit of the membrane, and the
entropy bill of every attempt to refuse it.

%=====================================================================
\section*{Acknowledgments}
The framework applied here is that of the companion papers
\cite{indivisible,instantaneous}; the debts to Wald and
Ishibashi--Wald, to the hyperboloidal program, and to the
near-horizon extension literature are recorded in
Section~\ref{subsec:notclaimed}.

%=====================================================================
\begin{thebibliography}{99}\small

\bibitem{indivisible}
J.~Torkelson,
\emph{Indivisible stochastic processes for non-Hermitian and
higher-derivative quantum systems: detailed balance, broken $\PT$
symmetry, and the explosion of ghosts},
companion paper (2026),
\url{https://github.com/embeddetech/p01-indivisible-ghosts}.

\bibitem{instantaneous}
J.~Torkelson,
\emph{Ghost vacuum decay as a stochastic process: instantaneous
states, the boost-orbit measure, and the foliation that ghosts make
visible},
companion paper (2026),
\url{https://github.com/embeddetech/p03-instantaneous-ghosts}.

\bibitem{Wald1980}
R.~M.~Wald,
\emph{Dynamics in nonglobally hyperbolic, static space-times},
J.\ Math.\ Phys.\ \textbf{21}, 2802 (1980).

\bibitem{IshibashiWald2003}
A.~Ishibashi and R.~M.~Wald,
\emph{Dynamics in non-globally-hyperbolic static space-times II:
general analysis of prescriptions for dynamics},
Class.\ Quantum Grav.\ \textbf{20}, 3815 (2003),
arXiv:gr-qc/0305012.

\bibitem{IshibashiWald2004}
A.~Ishibashi and R.~M.~Wald,
\emph{Dynamics in non-globally-hyperbolic static space-times III:
anti-de Sitter space-time},
Class.\ Quantum Grav.\ \textbf{21}, 2981 (2004),
arXiv:hep-th/0402184.

\bibitem{HorowitzMarolf1995}
G.~T.~Horowitz and D.~Marolf,
\emph{Quantum probes of spacetime singularities},
Phys.\ Rev.\ D \textbf{52}, 5670 (1995), arXiv:gr-qc/9504028.

\bibitem{ReggeWheeler1957}
T.~Regge and J.~A.~Wheeler,
\emph{Stability of a Schwarzschild singularity},
Phys.\ Rev.\ \textbf{108}, 1063 (1957).

\bibitem{MartelPoisson2001}
K.~Martel and E.~Poisson,
\emph{Regular coordinate systems for Schwarzschild and other
spherical spacetimes},
Am.\ J.\ Phys.\ \textbf{69}, 476 (2001), arXiv:gr-qc/0001069.

\bibitem{ReedSimon2}
M.~Reed and B.~Simon,
\emph{Methods of Modern Mathematical Physics II: Fourier Analysis,
Self-Adjointness},
Academic Press, New York (1975).

\bibitem{Feller1952}
W.~Feller,
\emph{The parabolic differential equations and the associated
semi-groups of transformations},
Ann.\ of Math.\ \textbf{55}, 468 (1952).

\bibitem{ItoMcKean}
K.~It\^o and H.~P.~McKean, Jr.,
\emph{Diffusion Processes and their Sample Paths},
Springer, Berlin (1965).

\bibitem{KarlinTaylor}
S.~Karlin and H.~M.~Taylor,
\emph{A Second Course in Stochastic Processes},
Academic Press, New York (1981).

\bibitem{Davies1980}
E.~B.~Davies,
\emph{One-Parameter Semigroups},
Academic Press, London (1980).

\bibitem{SzNagyFoias}
B.~Sz.-Nagy and C.~Foia\c{s},
\emph{Harmonic Analysis of Operators on Hilbert Space},
North-Holland, Amsterdam (1970).

\bibitem{Warnick2015}
C.~M.~Warnick,
\emph{On quasinormal modes of asymptotically anti-de Sitter black
holes},
Commun.\ Math.\ Phys.\ \textbf{333}, 959 (2015), arXiv:1306.5760.

\bibitem{GajicWarnick2020}
D.~Gajic and C.~M.~Warnick,
\emph{A model problem for quasinormal ringdown of asymptotically flat
or extremal black holes},
J.\ Math.\ Phys.\ \textbf{61}, 102501 (2020), arXiv:1910.08481.

\bibitem{Jaramillo2021}
J.~L.~Jaramillo, R.~Panosso Macedo, and L.~Al~Sheikh,
\emph{Pseudospectrum and black hole quasinormal mode (in)stability},
Phys.\ Rev.\ X \textbf{11}, 031003 (2021), arXiv:2004.06434.

\bibitem{Balachandran2018}
A.~P.~Balachandran, A.~R.~de~Queiroz, and A.~Saa,
\emph{Near-horizon modes and self-adjoint extensions of the
Schr\"odinger operator},
in \emph{Classical and Quantum Physics}, Springer Proc.\ Phys.\
\textbf{229} (2019), arXiv:1810.00020.

\bibitem{GuptaSen2002}
K.~S.~Gupta and S.~Sen,
\emph{Further evidence for the conformal structure of a Schwarzschild
black hole in an algebraic approach},
Phys.\ Lett.\ B \textbf{526}, 121 (2002), arXiv:hep-th/0112041.

\bibitem{CamblongOrdonez2005}
H.~E.~Camblong and C.~R.~Ord\'o\~nez,
\emph{Black hole thermodynamics from near-horizon conformal quantum
mechanics},
Phys.\ Rev.\ D \textbf{71}, 104029 (2005), arXiv:hep-th/0411008.

\bibitem{CamblongOrdonez2007}
H.~E.~Camblong and C.~R.~Ord\'o\~nez,
\emph{Conformal enhancement of holographic scaling in black hole
thermodynamics: a near-horizon heat-kernel framework},
JHEP \textbf{0712}, 099 (2007), arXiv:0709.2942.

\bibitem{FranchiLeJan2007}
J.~Franchi and Y.~Le~Jan,
\emph{Relativistic diffusions and Schwarzschild geometry},
Comm.\ Pure Appl.\ Math.\ \textbf{60}, 187 (2007),
arXiv:math/0410485.

\bibitem{Damour1978}
T.~Damour,
\emph{Black-hole eddy currents},
Phys.\ Rev.\ D \textbf{18}, 3598 (1978).

\bibitem{ThornePriceMacdonald}
K.~S.~Thorne, R.~H.~Price, and D.~A.~Macdonald (eds.),
\emph{Black Holes: The Membrane Paradigm},
Yale University Press, New Haven (1986).

\bibitem{ParikhWilczek1998}
M.~Parikh and F.~Wilczek,
\emph{An action for black hole membranes},
Phys.\ Rev.\ D \textbf{58}, 064011 (1998), arXiv:gr-qc/9712077.

\bibitem{STU1993}
L.~Susskind, L.~Thorlacius, and J.~Uglum,
\emph{The stretched horizon and black hole complementarity},
Phys.\ Rev.\ D \textbf{48}, 3743 (1993), arXiv:hep-th/9306069.

\bibitem{tHooftBrickWall1985}
G.~'t~Hooft,
\emph{On the quantum structure of a black hole},
Nucl.\ Phys.\ B \textbf{256}, 727 (1985).

\bibitem{tHooftAntipodal2016}
G.~'t~Hooft,
\emph{Black hole unitarity and antipodal entanglement},
Found.\ Phys.\ \textbf{46}, 1185 (2016), arXiv:1601.03447.

\bibitem{AMPS2013}
A.~Almheiri, D.~Marolf, J.~Polchinski, and J.~Sully,
\emph{Black holes: complementarity or firewalls?},
JHEP \textbf{02} (2013) 062, arXiv:1207.3123.

\bibitem{HorowitzMaldacena}
G.~T.~Horowitz and J.~Maldacena,
\emph{The black hole final state},
JHEP \textbf{02} (2004) 008, arXiv:hep-th/0310281.

\bibitem{KayWald1991}
B.~S.~Kay and R.~M.~Wald,
\emph{Theorems on the uniqueness and thermal properties of
stationary, nonsingular, quasifree states on space-times with a
bifurcate Killing horizon},
Phys.\ Rep.\ \textbf{207}, 49 (1991).

\bibitem{Unruh1981}
W.~G.~Unruh,
\emph{Experimental black hole evaporation?},
Phys.\ Rev.\ Lett.\ \textbf{46}, 1351 (1981).

\bibitem{BLV2005}
C.~Barcel\'o, S.~Liberati, and M.~Visser,
\emph{Analogue gravity},
Living Rev.\ Relativ.\ \textbf{29}, 1 (2026), arXiv:gr-qc/0505065.

\bibitem{CoutantParentani2012}
A.~Coutant, R.~Parentani, and S.~Finazzi,
\emph{Black hole radiation with short distance dispersion, an
analytical S-matrix approach},
Phys.\ Rev.\ D \textbf{85}, 024021 (2012), arXiv:1108.1821.

\end{thebibliography}

\end{document}
